%!TEX root = ./poly.tex

\section{Notation}
\label{sec:notation}

Here we collect all the notation used in the paper (as well as some notation which does not appear in this paper but is used in related works by the author). In fact, the notation system we introduce should be regarded as the ``fundamental technical lemma'' from which all else in the paper follows. (This was also largely the case in \cite{SulSlopes}.) We will also chronicle the evolution of the notation.

We use $=$ for canonical isomorphisms, and $\cong$ for non-canonical isomorphisms; the meaning of ``canonical'' is elaborated in Notation \ref{not:canonical}. We also write $x\approx y$ to mean $x=uy$ for a unit $u$.

The ring we use are
\begin{align*}
  R &= \text{general perfectoid ring}\\
  R^+ &= \text{$p$-torsionfree perfectoid ring}\\
  k &= \text{perfect $\F_p$-algebra}\\
  C &= \text{spherically complete perfectoid field}\\
  S &=
    \begin{cases}
      \text{quasiregular semiperfectoid ring} &\text{in \sec\ref{sec:qrsp}}\\
      \text{$p$-completely smooth $R$-algebra} & \text{in \sec\ref{sec:dRW}}
    \end{cases}\\
  A &=
  \begin{cases}
    \text{base $\delta$-ring} & \text{in \sec\ref{sec:qrsp}}\\
    \Ainf(R) & \text{in \sec\ref{sec:dRW}}
  \end{cases}\\
  W &= \rmW(k)
\end{align*}
We index $p$-typical Witt vectors and $p$-typical $\TR$ starting from $0$, so that $\TR^n = \THH^{C_{p^n}}$ and $\rmW_n(\F_p)=\Z/p^{n+1}$. We write $\bigTR$ for big $\TR$ and $\W$ for big Witt vectors. We write $\phi$ for the Frobenius on $A$, $\varphi$ for the Frobenius on prismatic cohomology, and $\~\varphi$ for a divided Frobenius.

If $(A, I)$ is a prism, then $I_n = I\phi(I)\dotsm\phi^{n-1}(I)$ \cite[Notation 2.2.2]{APC}. We recall the following elements from \cite[Lemma 3.12]{BMS1}: if $R$ is a perfectoid ring with corresponding perfect prism $(A, I)$, then
\begin{align*}
  \~\xi & \text{ is a generator of }I\\
  \xi &= \phi^{-1}(\~\xi) \text{ is a generator of }\Nyg^{\ge 1}\prism_R = \phi^{-1}(I)\\
  \~\xi_n &= \~\xi\phi(\~\xi)\dotsm\phi^{n-1}(\~\xi) \text{ is a generator of } I_n\\
  \xi_n &= \xi\phi^{-1}(\xi)\dotsm\phi^{-(n-1)}(\xi) \text{ is a generator of } \phi^{-n}(I_n)
\end{align*}
Beware that this notation is off-by-one from that used in \cite{PrismsTambara1}.

In the $q$-de Rham case, we have the element $q=[\epsilon]\in\Ainf(R)$ from from \cite[Proposition 3.17]{BMS1}, allowing us to make these more concrete as $q$-analogues:
\begin{align*}
  \mu &= 1-q\\
  \~\xi_n &= (p^n)_q = \frac{\phi^n(\mu)}\mu\\
  \xi_n &= (p^n)_{q^{1/p^n}} = \frac{\mu}{\phi^{-n}(\mu)}
\end{align*}

We will sometimes use the \emph{relative} theory
\[ \TR(S \mid \SW k) \quad\text{or}\quad \TR(S \mid \SAinf{R}) \]
where $\SAinf R := \SW{R^\flat}$, and $\SW k$ is the functor of spherical Witt vectors \cite[Example 5.2.7]{Ell2}, \cite{AntieauSWitt}, \cite[\sec2]{ChromaticNullstellensatz}, \cite{AlternativeSWitt}. This is equivalent to $\TR(S;\Z_p)$ by \cite[Lemma 5.20]{Zhouhang}, an avatar of the identification between absolute and relative prismatic cohomology under a perfect prism \cite[Lemma 4.8]{Prismatic}, \cite[Proposition 4.4.12]{APC}, \cite[Remark 11.9]{BMS2}. However, the relative notation is helpful for emphasizing the base in results that depend on it, and for reducing the context needed when reading our formulas on the go.

The circle group is denoted by $\T$. We write $\lambda^n$ for the 1-dimensional complex $\T$-representation in which $z\cdot w=z^n w$ for $z\in\T$ and $w\in\C$; this is the $n^{\text{th}}$ power of $\lambda=\lambda^1$ in the complex representation ring $\RU(\T)$. We set $\lambda_i = \lambda^{p^i}$, with $\lambda_\infty = \lambda^0 = \C^\triv$. Every real representation of $\T$, and every real orientable representation of $C_n$, is of the form $\alpha$ or $\alpha-1$ for a complex representation $\alpha$. %We do not consider sign representations (when $p=2$) in this paper.

\begin{notation}[Dimensional encoding]
  The first major maneuver is to index a representation $\alpha$ by its (eventually constant) dimension sequence $d_\bullet(\alpha)=\dim_\C \alpha^{C_{p^\bullet}}$, rather than the coefficients of its irreducible decomposition. This insight is due to Angeltveit-Gerhardt \cite{ROS1TR}, and would be a good lesson for the rest of equivariant homotopy theory to learn---for example, it changes the ``double cone'' pattern seen in computations into a simple quadrant pattern.

    To appreciate the difference, note that the dimensional encoding is suited for computing with the isotropy separation square, while the irreducible encoding is suited for computing with cell structures. The isotropy separation method is multiplicative, and this is a major advantage over the cellular method, which is not. See \cite{Zengv1} for a detailed comparison of these methods for computing $\mpi_\rog(\Z)$ over $C_{p^2}$.

    Thus, we write a virtual representation $\alpha$ as $\alpha=(d_0,\dotsc,d_n;d_\infty)$, where $d_i = \dim_\C(\alpha^{C_{p^i}})$, and $d_\infty$ is the tail of the eventually constant sequence. Note that $(d_0,\dotsc,d_n;d_\infty) = (d_0,\dotsc,d_n,d_\infty;d_\infty)$; these two encodings give slightly different spectral sequences, but the end result is of course the same. If $d_i(\alpha)=d_\infty$ for $i\ge n$, we sometimes refer to $\alpha$ as a $C_{p^n}$-representation, although this is slightly abusive.

  To convert between this encoding and the irreducible decomposition $\alpha=k_0\lambda_0 + \dotsb + k_n \lambda_n + k_\infty \lambda_\infty$, the formulas are
  \begin{align*}
    d_i &= \sum_{i\le j\le \infty} k_j\\
    k_i &= d_i - d_{i+1} \qquad i <\infty\\
    k_\infty &= d_\infty
  \end{align*}

  Beware that the above encoding does not distinguish between $\lambda^n$ and $\lambda^{p^{v_p(n)}}$. This is harmless for most purposes, but it does twist restriction ($F$) maps by the unit $n/p^{v_p(n)}$, and one needs to keep track of this when trying to pin down the Norm exactly (rather than just up to units).
\end{notation}

\begin{notation}[HOTRSS]
  The next powerful move is to unravel the induction computing $\TR^n_\rog$ into a spectral sequence, the homotopy orbits to TR spectral sequence. This again is due to Angeltveit-Gerhardt \cite[\sec3]{ROS1TR}. We gave a more conceptual construction of this spectral sequence (that makes it easy to see the multiplicative and Mackey structure) in \cite[\sec3]{SulROGTHH}, and it is this formulation that we will use. We will always write this as a two-column spectral sequence in Adams grading, where the left column computes $\pi_{\alpha-1}$ and the right column computes $\pi_\alpha$.
\end{notation}

\begin{notation}[Fool's gold]
  To make the multiplicative and Mackey structures transparent, we shift to indexing using the gold elements. We learned this from Dylan Wilson and \cite{Zeng}. Actually, one of the major advances of this paper is to \emph{remove} the gold elements---but first let us just recall what they are.

  Given an actual representation $V$, we write $a_V \in \pi^G_{-V} \S_G$ for the \emph{Euler class} of $V$, obtained by suspending the inclusion $\{0\}\mono V$ (see \cite[Definition 3.11]{HHR}). In particular we can map $a_V$ to any $G$-spectrum $X$. Note that this manifestly does not depend on any choices, i.e.\ it is ``canonical''. We will not actually get rid of the $a_V$, but they will be hidden in what follows.

  In favorable cases there are also \emph{Thom classes} $u_V \in \pi^G_{V-\dim_\R V}(X)$ giving an isomorphism
\[ u_V\colon \Sigma^V X \iso \Sigma^{\dim_\R V} X. \]
Such an isomorphism typically only exists when $X$ is Borel-complete; however, if $X$ is a genuine $G$-spectrum and $X^h=X^{EG_+}$ is its Borel completion, the classes $u_V$ (but not $u_V^{-1}$) often lift along $X\to X^h$. For the case we need, the $u_V$ are constructed in \cite[\sec4]{RiggenbachDoublePoints}. Other good sources to learn about these classes are \cite[\sec5.2]{Zeng} and \cite{HZhZ}. Crucially, the $u_V$ are \emph{non-canonical}, so these are the ones we want to get rid of.

  The $n^\th$ power maps $S^{\lambda^k} \to S^{\lambda^{nk}}$ give divisibility relations $a_{\lambda^k} \mid a_{\lambda^{nk}}$, and we set
  \begin{align*}
    a_i &\defeq a_{\lambda_{i-1}}^{-1} a_{\lambda_i}\\
    u_i &\defeq u_{\lambda_{i-1}} u_{\lambda_i}^{-1}
  \end{align*}
  When $i=0$, we interpret these formulas by setting $a_{\lambda_{-1}} \defeq 1$ and $u_{\lambda_{-1}} = \sigma$, the B\"okstedt generator. (There is some interesting philosophy here). Setting $T=\THH(R;\Z_p)$, these classes are summarized in the following diagram:
  \[\xymatrix@=5em{
    & T \ar@/_.5em/[dl]_-{a_{\lambda_{-1}} \defeq \id} \ar[d]^{a_{\lambda_0}} \ar[dr]^{a_{\lambda_1}} \ar@/^1em/[drr]^-{a_{\lambda_2}}\\
      T \ar@<.5ex>[r]^-{a_0} & \Sigma^{\lambda_0} T \ar@<.5ex>[l]^-{u_0} \ar@<.5ex>[r]^-{a_1} & \Sigma^{\lambda_1} T \ar@<.5ex>[l]^-{u_1} \ar@<.5ex>[r]^-{a_2} & \cdots \ar@<.5ex>[l]^-{u_2}\\
                            & \Sigma^2 T \ar@/^.5em/[ul]^-{u_{\lambda_{-1}} \defeq \sigma} \ar[u]_-{u_{\lambda_0}} \ar[ur]_-{u_{\lambda_1}} \ar@/_1em/[urr]_-{u_{\lambda_2}}
  }\]
  Note that the generator $t$ of $\TP_{-2}(R;\Z_p)$ is given by $t=a_{\lambda_0} u_{\lambda_0}^{-1}$.

  The gold classes are related by the ``$q$-gold relations''
  \begin{alignat*}{2}
    \phi^{-1}(\~\xi_{i+1})a_{\lambda_i} &= 0\\
    a_{\lambda_j} u_{\lambda_i} &\approx \phi^i(\~\xi_{j-i}) a_{\lambda_i} u_{\lambda_j} &\quad& (-1\le i<j)\\
    a_i u_i &\approx \phi^i(\xi)\\
    \intertext{from \cite[Lemma 4.9]{SulSliceTHH}, as well as the ``infinite products''}
    u_{\lambda_{i-1}} &\approx u_i u_{i+1} u_{i+2} \dotsm
  \end{alignat*}
  These are up to units since the $u_V$ themselves are only defined up to units. The infinite products can be made precise by viewing $\TF(R;\Z_p) = \lim_F \TR^n(R;\Z_p)$ as a condensed spectrum with the limit topology, or by using Breuil-Kisin twists.

  In \cite{SulROGTHH}, we used the classes
  \[ \vartheta^\alpha_r \defeq \left(\prod_{i=0}^r a_i^{-d_i}\right)\left(\prod_{i=r+1}^\infty u_i^{d_i}\right) \]
  using the above convention for infinite products. More canonically (Notation \ref{not:canonical}), this corresponds to tensoring with the ideal $\bigotimes_{i=r}^\infty \phi^{i}(I)$.
\end{notation}

\begin{remark}[Polar rings]
  It is often said that the multiplication on $\RO(G)$ is a red herring, and that the notion of $\RO(G)$-graded homotopy only sees the abelian group structure of $\RO(G)$. Above, we noticed that this is not entirely true: $n^\th$ power maps give divisibility relations $a_V \mid a_{V^{\tnsr n}}$.

  Bauer \cite{BauerPolar} has defined the notion of a \emph{$p$-polar ring}: an abelian group equipped with an operation ``$x\mapsto x^p$'', but not a full multiplication; in particular, he notes that $p$-typical Witt vectors are naturally defined for such. One could remove the $p$-typical restriction and define \emph{polar rings} as abelian groups equipped with operations ``$x\mapsto x^n$'', and then say that $\RO(G)$-graded homotopy sees $\RO(G)$ as a polar ring.
\end{remark}

\begin{conjecture}[$G$-polar derivations and the $q$-de Rham complex]
  More generally, one could define \emph{naive polar $G$-rings} as modules with $G$-action equipped with a notion of ``$x \mapsto x^{G/H}$'', and \emph{genuine polar $G$-rings} as Mackey functors equipped with norm maps but not a full multiplication. Here is a reason for doing so: Hill \cite{HillTambaraAQ} has defined a cotangent complex for Tambara functors, based on a notion of $G$-derivation. A $G$-derivation is a levelwise derivation which additionally satisfies compatibility with Norm maps, which we state heuristically as
  \[ d(x^{G/H}) = G/H \cdot \text{``$x^{G/H} \dlog x$''}. \]

  We can turn the $q$-de Rham complex into a complex of Mackey functors via $C_n \mapsto \frac{\qOmega^*}{(q^n-1)}$ or $C_n \mapsto \frac{\qOmega^*}{(n)_q}$, and we can turn $\qOmega^0$ into a Tambara functor using $q$-twisted powers \cite[\sec5.3]{PrismsTambara1}. The $q$-derivative then satisfies the above compatibility with norms, but it is not a levelwise derivation. We conjecture that, suitably defined, the cotangent complex of $G$-polar rings (which should remove the derivation condition, but keep the norm condition) recovers $\qOmega^1$; here the group is $G=\Q/\Z$, and perhaps the cotangent complex should be relative to the Tambara functor $\m\RU = \left[C_n \mapsto \frac{\Z[q]}{q^n-1}\right]$.
\end{conjecture}

\begin{notation}[Poly-representations]
  To study $\TR_\rog$ rather than just $\TR^n_\rog$, one switches to using poly-representations, introduced in \cite[\sec4.2]{SulSlopes} in order to conceptualize \cite[Proposition 4.2.3]{HMCyclicPolytopes}.

  In the $\TR$ setting, we denote a poly-representation $\m\alpha$ by $(\m d_0, \m d_1, \m d_2,\dotsc)$, with $\m d_n = \dim_C \m\alpha(C_{p^n})$. We have poly-graded gold classes
  \begin{align*}
    \m a_{\lambda_i} \colon C_{p^n} &\mapsto a_{\lambda_{n-i}}\\
    \m u_{\lambda_i} \colon C_{p^n} &\mapsto u_{\lambda_{n-i}}\\
    \m a_i \colon C_{p^n} &\mapsto a_{n-i}\\
    \m u_i \colon C_{p^n} &\mapsto u_{n-i}
  \end{align*}
  The polygonic analogue of $\vartheta^\alpha_r$ is
  \[ \vartheta^{\m\alpha}_r = \m u_0^{\m d_0}\u_1^{\m d_1}\dotsm\u_r^{\m d_r} \a_{r+1}^{-\m d_{r+1}}\dotsm \]
  which more canonically corresponds to tensoring with $\bigotimes_{i=0}^r \phi^{-i}(\xi)^{\m d_i}A$ (Definition \ref{def:poly-ideal}).
\end{notation}

\begin{warning}
  We should really write $d_{-n}(\m\alpha)$, but this notation gets too bulky; the underline notation emphasizes the connection to Mackey functors, and $\m d_n$ still looks pretty close to $d_{-n}$. However, notational macros like this can obscure the bigger picture: for integral poly-representations, one should write $d_{1/n}(\m\alpha)$.
\end{warning}

In particular, interpreting $\m a_i = a_{-i}$ and $\m u_i = u_{-i}$, the polygonic gold relation is still just $a_i u_i = \phi^i(\xi)$, but for negative $i$. In $\TR_\rog(\O_C;\Z_p)$, we have the additional relation
\[
  \beta a_{\m\lambda_i} = \phi^{-i}(\mu) u_{\m\lambda_i}
\]
where $\beta$ is the Bott element.

We highlight some significant families of poly-representations; these are not used in the current paper, but are important for understanding and used in related works.

\begin{definition}
  The \emph{irreducible} poly-representations are
  \[ \lambda^{1/k}\colon C_n \mapsto \begin{cases} \lambda^{n/k} & k\mid n\\ 0 & k\nmid n\end{cases} \]
  In the $p$-typical case, we set $\m\lambda_i = \lambda^{1/p^i}$. We have $\m d_j(\m\lambda_i) = [j \ge i]$, where $[P]$ is the Iverson bracket.
\end{definition}

\begin{definition}
  \label{def:regular-poly}
  For a positive rational number $\nu\in\Q_{>0}$, the \emph{regular poly-representation of slope $\nu$} is
  \[ \hsl\nu\colon C_n \mapsto [\lc n\nu\rc]_\lambda \defeq \lambda^0 + \dotsb + \lambda^{\lc n\nu\rc-1}. \]
  These are so-named because $\hsl1$ specializes to the complex regular representation of each $C_n$. Beware that the irreducible decomposition of $\hsl\nu$ is infinite. This has dimension sequence $\m d_i(\hsl\nu) = \lc p^i\nu\rc$, or integrally $d_{1/n}(\hsl\nu)=\lc n\nu\rc$.
\end{definition}

The notational evolutions reviewed thus far have already appeared in previous papers. Now we explain the new maneuvers introduced in this paper.

\begin{notation}
  To extend the homotopy orbits to TR spectral sequence from perfectoid rings to smooth algebras over such, the main simplification is to pull out
  \[ \TR^\bullet_{\rog,j}(S;\Z_p) \defeq \text{contribution of $\dRW[\bullet]^j_S$ to $\TR^\bullet_{\rog+j}(S;\Z_p)$.} \]
  That we can do this follows by [the proof of] Hesselholt's theorem \cite{LarsPTyp} (see also \cite[Theorem 11]{HesselholtHandbook})
  \[ \TR^n_*(S;\Z_p) =  \dRW[n]^*_{S/R} \otimes_{\rmW(R)} \TR^n_*(R;\Z_p), \]
  which in this notation becomes
  \[ \TR^n_{2i,j}(S;\Z_p) = \dRW[n]^j_{S/R} \tnsr \sigma_n^i \]
  where $\sigma_n$ is the B\"okstedt generator of $\TR^n_2(R;\Z_p)$. With this notation, the computation for a smooth algebra $S$ works mostly the same as for the base perfectoid ring $R$.
\end{notation}

\begin{notation}[Canonical identifications]
  \label{not:canonical}
  Since the Thom classes $u_{\lambda^i}$ are invertible in $\TCmin$, we have \emph{non-canonical} isomorphisms
  \begin{align*}
    \TCmin_\alpha(S;\Z_p) &\cong \TCmin_{2 d_0(\alpha)}(S;\Z_p)\\
    \TP_\alpha(S;\Z_p) &\cong \TP_{2 d_0(\alpha)}(S;\Z_p)\\
                      &\cong \TP_0(S;\Z_p)
  \end{align*}
  for every $\alpha\in\RU(\T)$. The content of our computation is to provide \emph{canonical} identifications. Since it is one of the main focuses of this paper, and the author was confused about it for several years, let us clarify what ``canonical'' really means.

  Usually we are told it means ``without making choices'', but this gets a bit philosophical. In particular, we want to say that ``$p\Z_p$ is isomorphic to $\Z_p$, but not canonically''. One could try to justify this by arguing that there are $\Z_p^\times$ possible generators, but this falls apart on noting that $p\Z_p = p\N \tnsr_\N \Z_p$ and $|\N^\times|=1$.

  On the other hand, $p\Z_p$ has a natural map to $\Z_p$, and $(p\Z_p \to \Z_p)$ is certainly not isomorphic to $(\Z_p \morph^\id \Z_p)$ as objects of the slice category $\Mod_{/\Z_p}$. In general, let $T$ be a $G$-spectrum, and write a virtual $G$-representation $\alpha$ as $\alpha = i + U - V$, where $U$ and $V$ are actual $G$-representations with $U^G=V^G=0$, and $i\in\Z$. Then the equviariant Euler classes give $\pi_\alpha^G T$ the structure\footnote{Strictly speaking, this depends on the choice of $U$ and $V$; there is an instructive sermon about this in \cite[\sec6]{Adams}, see also \cite{DuggerCoherence,GradedTambara} for a careful discussion. We avoid this issue by fixing the basis $\{\lambda^i\}$ of $\RU(\T)$ and never sharing irreducibles between $U$ and $V$.} of an object in $\mathrm{Span}(\mathrm{Mod})_{\pi^G_i(T)/}$:
  \[ \xymatrix{
    & \ar[dl]_-{a_U} \pi^G_{i+U} T \ar[dr]^-{a_V}\\
    \pi_i^G T && \pi^G_{i+U-V}T
  } \]
  Now by ``canonical'' identifications of the $\RO(G)$-graded groups, at minimum we mean we are going to identify the maps in this span. There are lots more maps coming from the equivariant structure ($F,\,V,\,N,\,R,\,\dots$), and we want to identify these as well.

  We do this not by writing down a bunch of matrices, but by expressing the answer as tensored with various line bundles---a sort of dimensional analysis---so that all the maps are simply ``the obvious ones''. For example, in the $\TR$ case there are $F$ and $R$ maps, which act on the B\"okstedt generator by
  \begin{align*}
    F(\sigma_n) &\approx \sigma_{n-1}\\
    R(\sigma_n) &\approx \phi^{-n}(\xi)\sigma_{n-1}
  \end{align*}
  These can be explained by the natural maps
  \[\xymatrix{
      \xi_{n+1} A \ar@<1ex>[r]^-F \ar@<-1ex>[r]_-R & \xi_n A
  }\]
  where $R$ is the inclusion, and $F(f)=\phi(f)/\~\xi$ is a divided Frobenius. Thus, we can further rewrite Hesselholt's theorem as
  \[ \TR^n_{2i,j}(S;\Z_p) = \dRW[n]^j_{S/R} \tnsr_A \xi_{n+1}^i A \]
  Actually, the formulas for $F$ and $R$ only justify this once we know the base case $\THH_{2i,j}(S;\Z_p) = \Omega^j_{S/R} \tnsr_A \xi^j A$. This follows from the comparison with the cotangent complex, as explained in Example \ref{ex:tor-ext}.

  The module $\xi_nA/\xi_n^2A$ is often denoted by $(A/\xi_n)\{1\}$, but we will need a much larger range of twists than this, so we will instead write this as $(A\tnsr \xi_n)/\xi_n$ or $(A/\xi_n) \tnsr \xi_n$. In general, if $M$ is an $A$-module, a notation like $M \tnsr \xi^2 \phi^{-1}(\xi)^{-3}$ means
  \[ M \tnsr_A \xi^2A \tnsr_A \Hom_A\left(\phi^{-1}(\xi)^3A,A\right) \]

We will sometimes need to make sense of this notation for certain \emph{infinite} products. In the prismatic context, the infinite product $I\phi(I)\dotsm$ can be made sense of as the Breuil-Kisin twist, as in \cite[\sec2.2]{APC}. In the de Rham-Witt context, the symbol $\xi\phi^{-1}(\xi)\dotsm$ means the ideal
  \[ \bigcap_{n=0}^\infty \xi_n \Ainf. \]
  We only consider this intersection when the perfectoid base is spherically complete, in which case this intersection is equal to $(1-q)\Ainf$ by \cite[Lemma 3.23]{BMS1}. We do not know what the intersection looks like in any other cases (besides $\phi^{-r}(\xi)\phi^{-(r+1)}(\xi)\dotsm = \phi^{-r}(1-q)$).

  When $R$ is a perfectoid ring, the canonical identification $\TF_{\lambda}(R;\Z_p) = \xi\Ainf(R)$ was explained to the author by Mike Hill at MSRI in 2019. Unreadable canonical identifications of $\TF_\rog(R;\Z_p)$ and $\TR^\bullet_\rog(R;\Z_p)$ in the positive and negative ranges appeared in the author's PhD thesis. We discovered the $q$-gold relation \cite[Lemma 4.9]{SulSliceTHH} a few months later; this essentially implies the canonical identifications, but it took years for the author to understand this. The first canonical identifications in the generality of quasiregular semiperfectoid rings are due to Riggenbach \cite{RiggenbachTruncated}.
\end{notation}

\begin{example}
  \label{ex:tor-ext}
  Let $A$ be a ring, $f\in A$ a non-zerodivisor, $I=(f)$ the principal ideal generated by $A$, and $\bar A=A/I$. It is well-known that for any $A$-module $M$,
  \begin{align*}
    \Tor^A_1(\bar A, M) &\cong {}_I M \fade{{}\defeq \{m \in M \mid Im = 0\}}\\
    \Ext^1_A(\bar A, M) &\cong M/I
  \end{align*}
  as one sees by resolving $A/I$ as $A\morph^f A$.

  It is less well-known that these isomorphisms are non-canonical: the chosen resolution depends on $f$, while $\bar A$ only depends on $I$. Let $\bar A\{1\} = I/I^2$, and for any $\bar A$-module $N$ let $N\{n\} = N \tnsr_{\bar A}(\bar A\{1\})^{\tnsr n}$. In general, we have
  \begin{align*}
    \gr^i_I A &= \begin{cases}\bar A\{i\} & i\ge 0\\0 & i <0\end{cases} &
    \gr^i_I (A[1/I]) &= \bar A\{i\}
  \end{align*}

  Now a \emph{canonical} free resolution of $\bar A$ is $I\to A$, which gives \emph{canonical} identifications
  \begin{align*}
    \Tor^A_1(\bar A, M) &= {}_I M \{1\}\\
    \Ext^1_A(\bar A, M) &= M/I\{-1\}.
    \intertext{Similarly, there are canonical identifications}
    \mathrm L_{\bar A\mid A} &= \bar A\{1\}[1]\\
    \HH_{2i}(\bar A \rel A) &= \Gamma^i \bar A\{i\}.
  \end{align*}
  In particular, if there is additional structure like a Galois action or a Frobenius floating around, these twists will keep track of the correct actions.

  Please mention this point the next time you teach homological algebra.
\end{example}


\begin{notation}[The power of negative thinking] In both \cite{SulROGTHH} and \cite{SulSlopes}, there is a lot of fussing over whether a given expression is negative or not. Usually this is because we have some filtration like $\xi^i A$, and when $i$ is negative we want to interpret this as $\xi^{\max(i, 0)}A$ (i.e.\ make the filtration constant in negative degrees), rather than the $\xi$-adic filtration on the localization $A[\xi^\pm]$. By interpreting filtrations this way, and similarly defining $\rmW_i(R)=0$ for $i<0$, we can avoid handling $2^\infty-1$ cases in exchange for complicated expressions like $\dRW[-(\m d_1+\m d_2)]^j$ that will simplify in all practical situations.
\end{notation}

We close with a final maneuver which we have not adopted, but could be useful in future.

\begin{notation}[Accepting what we cannot know] In \sec\ref{sec:dRW}, one reason for restricting to the contiguous range is that we do not know how to compute the kernels of maps such as
  \[ \frac A\xi \xra{\phi^{-1}(\xi)^a} \frac A{\xi\phi^{-1}(\xi),\, \phi^{-2}(\xi)^b} \]
  where $a,b\ge0$ are arbitrary. But we could simply give this kernel a name, and compute other things in terms of it, and leave it to the commutative algebraists to give a more explicit description of its structure.
\end{notation}
